\subsection{POWERS IN STANDARD GROUPS}

We preserve the notation of no. 4.

PROPOSITION 8. Let \(n \in \mathbf{Z}\) and \(h_n\) be the mapping \(x \mapsto x^n\) of \(G\) into \(G\) . Let \(a\) be a non-zero ideal of \(A\) contained in \(m\) , such that \(n \notin a\) . Then \(h_n|G(a)\) is an isomorphism of the analytic manifold \(G(a)\) onto the analytic manifold \(G(na)\) .

By definition of standard groups, \(h_n\) is equal on the whole of \(G\) to the sum of an integral series with coefficients in \(A^r\) . By § 5, formula (4), this series is of the form

\[
h _ {n} (x) = n x + \sum_ {| \alpha | \geqslant 2} a _ {\alpha} x ^ {\alpha}.
\]

Hence, for \(x \in G\) ,

\[
\begin{array}{r l} h _ {n} (n x) & = n ^ {2} \Big (x + \sum_ {| \alpha | \geqslant 2} a _ {\alpha} n ^ {| \alpha | - 2} x ^ {\alpha} \Big) \\ & = n ^ {2} \mathrm{S} (x) \end{array}
\]

where we write \(\mathrm{S}(x)=x+\sum_{|\alpha| \geq 2} a_{\alpha} n^{|\alpha|-2} x^{\alpha}\) . This series \(\mathrm{S}(x)\) defines an analytic mapping, also denoted by S, of G into G. By Algebra, Chapter IV, §6, Proposition 8, there exists an integral series \(S'\) in r variables with coefficients in \(A^r\) such that \(\mathrm{S}'(\mathrm{S}(\mathrm{X}))=\mathrm{S}(\mathrm{S}'(\mathrm{X}))=\mathrm{X}\) . Hence S is an isomorphism of the analytic manifold G onto itself and, for every non-zero ideal b of A contained in m, \(\mathrm{S}(\mathrm{G}(\mathrm{b}))\subset\mathrm{G}(\mathrm{b}),\mathrm{S}'(\mathrm{G}(\mathrm{b}))\subset\mathrm{G}(\mathrm{b})\) and therefore

\[
\mathrm{S} (\mathrm{G} (\mathfrak {b})) = \mathrm{G} (\mathfrak {b}).
\]

As \(h_n(y) = n^2\mathrm{S}\left(\frac{1}{n} y\right)\) for \(y \in n\mathrm{G}\) , we see that \(h_n|n\mathrm{G}(\mathfrak{b})\) is an isomorphism of the analytic manifold \(n\mathrm{G}(\mathfrak{b})\) onto the analytic manifold \(n^2\mathrm{G}(\mathfrak{b})\) . But, as \(n \notin \mathfrak{a}\) , \(|n| > |\lambda|\) for all \(\lambda \in \mathfrak{a}\) , hence \(n^{-1}\mathfrak{a} \subset \mathfrak{m}\) and hence \(\mathfrak{a}\) is of the form \(n\mathfrak{b}\) where \(\mathfrak{b}\) is a non-zero ideal of \(\mathrm{A}\) contained in \(\mathfrak{m}\) .

COROLLARY. If \(n\) is invertible in A, \(h_n\) is an isomorphism of the analytic manifold G onto itself. For every non-zero ideal \(\mathfrak{a}\) of A contained in \(\mathfrak{m}\) , \(h_n(\mathrm{G}(\mathfrak{a})) = \mathrm{G}(\mathfrak{a})\) . For all \(x \in \mathrm{G}\) , \(\omega(x^n) = \omega(x)\) .

This follows immediately from Proposition 8.

PROPOSITION 9. Suppose that \(p \neq 0\) .

\begin{enumerate}
\def\labelenumi{(\roman{enumi})}
\item
  Let \(\mathfrak{a},\mathfrak{b}\) be non-zero ideals of \(\mathbf{A}\) such that \(\mathfrak{b} \subset \mathfrak{a} \subset \mathfrak{m}\) . In the group \(\mathrm{G}(\mathfrak{a}) / \mathrm{G}(\mathfrak{b})\) , every element has order a power of \(p\) .
\item
  Suppose that \(v(p) = 1\) . If \(x \in G\) is such that \(\omega(x) > \frac{1}{p - 1}\) , then
\end{enumerate}

\[
\omega (x ^ {p}) = \omega (x) + 1.
\]

By § 5, formula (4), for all \(x \in G\) ,

\[
x ^ {p} = p x + \sum_ {| \alpha | \geqslant 2} c _ {\alpha} x ^ {\alpha}
\]

where \(c_{\alpha} \in A^{r}\) for all \(\alpha\) . Even for proving (i) it can be assumed that \(v(p) = 1\) . Then if \(\omega(x) \geqslant 1\) , it follows that \(\omega(x^{p}) \geqslant \omega(x) + 1\) and hence \(\omega(x^{p^{n}})\) tends to \(+\infty\) as \(n\) tends to \(+\infty\) ; this proves (i). As \(\binom{p}{i}\) is divisible by \(p\) for

\(1 \leqslant i \leqslant p - 1\) , Proposition 2 of § 5, no. 3, proves that \(c_{\alpha} \in pA^{r}\) for

\[
2 \leqslant | \alpha | \leqslant p - 1
\]

and hence

\[
\omega \left(c _ {\alpha} x ^ {\alpha}\right) > \omega (p x) = \omega (x) + 1 \quad \text { for } 2 \leqslant | \alpha | \leqslant p - 1.
\]

On the other hand, if \(|\alpha| \geqslant p\) , \(\omega(c_{\alpha}x^{\alpha}) \geqslant p\omega(x)\) and \(p\omega(x) > \omega(x) + 1\) if \(\omega(x) > \frac{1}{p - 1}\) . This proves (ii).

